% !TeX spellcheck = es_ES

\documentclass{article}

\usepackage{amsmath}
\usepackage[spanish]{babel}
\usepackage{enumerate}
\usepackage{hyperref}
\hypersetup{
	colorlinks,
	citecolor=black,
	filecolor=black,
	linkcolor=black,
	urlcolor=black,
}
\usepackage{tcolorbox}
\tcbuselibrary{raster}

\newcommand{\sectionx}[1]{\section*{#1}\label{sec:#1}\addcontentsline{toc}{section}{\nameref{sec:#1}}}
\newcommand{\subsectionx}[1]{\subsection*{#1}\label{subsec:#1}\addcontentsline{toc}{subsection}{\nameref{subsec:#1}}}
\newcommand{\subsubsectionx}[1]{\subsubsection*{#1}\label{subsubsec:#1}\addcontentsline{toc}{subsubsection}{\nameref{subsubsec:#1}}}

\title{Apuntes de Introducción al Pensamiento científico CBC (40) \\
	Cátedra Campins, curso 62415\\
	1° C 2003}
\author{Darío Eduardo Ramos}

\begin{document}
	
\maketitle

\tableofcontents{}

\newpage

\sectionx{Resumen del libro de Alicia Gianella}

\subsectionx{I. Introducción: ¿qué es la epistemología?}

Dentro de las ciencias del conocimiento, se puede mencionar:

\begin{enumerate}[a.]
	\item Gnoseología: analiza aspectos del conocimiento desde la filosofía.
	\item \textbf{Epistemología}: se centra en el conocimiento científico. Para algunos autores, es ``la ciencia de la ciencia''. Una vertiente más moderna es la llamada ``epistemología naturalizada'', y propone un enfoque más interdisciplinario.
	\item Metodología de la ciencia: en un plano ya más práctico, lidia con técnicas de investigación.
	\item Filosofía de la ciencia: un poco más abarcativa que la epistemología pura: incluye ética, ontología y otras disciplinas.
\end{enumerate}

\subsectionx{II. Aspectos estructurales y dinámicos del conoc. científico}

La epistemología analiza , por un lado, aspectos estructurales como los componentes del conocimiento científico y sus interrelaciones. Por ejemplo, hipótesis y teorías científicas y cómo las primeras evolucionan hacia las segundas. Por otro lado, también analiza aspectos dinámicos como las creencias, sus cambios a lo largo del tiempo, y sus relaciones con otras disciplinas.

Según \textbf{Reichenbach}, hay 3 contextos epistemológicos:

\begin{enumerate}
	\item \textbf{Descubrimiento}: surgimiento, aceptación o rechazo de ideas nuevas. Existe un debate dentro de la epistemología respecto a si este contexto es un tema que le concierne: para Reichenbach y Popper, por ejemplo, la respuesta es no.
	\item \textbf{Justificación}: evaluación de hipótesis y teorías obviando sus circunstancias y contingencias (controversia).
	\item \textbf{Aplicación}: uso práctico de la ciencia en tecnología.
\end{enumerate}

Otro autor, Lakatos, propone un modelo alternativo:

\begin{enumerate}
	\item Historia externa: circunstancias, aceptación o rechazo.
	\item Historia interna: justificación empírica.
\end{enumerate}

\subsubsectionx{1. Componentes del conocimiento y sus relaciones}

Los componentes del conocimiento se denominan \textbf{saberes}. Los mismos pueden ser:

\begin{itemize}
	\item \textbf{Proposicionales}: se sabe algo verdadero, se cree y está fundamentado. En una concepción aristotélica, por \textbf{verdad} se entiende la correspondencia entre una proposición y un estado de cosas real. Y en cuanto al fundamento, sin el mismo se considera que no hay conocimiento.
	\item \textbf{Instrumentales}: Si los saberes proposicionales implican saber el qué, los instrumentales implican saber el cómo. Son el conocimiento enfocado en ejecutar acciones, procedimientos o habilidades.
\end{itemize}

\subsubsectionx{2. Enunciados y relaciones}

Los enunciados son oraciones \textbf{expresivas}, \textbf{directivas} e \textbf{informativas}. Esto último las hace \textbf{proposiciones}, porque afirman o niegan algo. Se les puede asignar un valor de verdad.

En cuanto a las \textbf{relaciones}, requieren:

\begin{itemize}
	\item \textbf{Consistencia}: que no haya contradicciones. Demostrar todo equivale a demostrar nada.
	\item \textbf{Equivalencia}: mismo valor de verdad.
	\item \textbf{Implicación}: si la premisa se cumple, la conclusión debe cumplirse.
	\item \textbf{Deducción}: si la premisa no se cumple, la conclusión no puede cumplirse.
\end{itemize}

\subsubsectionx{3. Razonamientos}

Un \textbf{razonamiento deductivo} es un conjunto de enunciados en los que uno, la \textbf{conclusión}, se desprende de los otros, las \textbf{premisas}.

\begin{tcolorbox}[colback=blue!5!blue!20!white,colframe=blue!75!black,title=Noción central en lógica, fontupper=\large]
Un razonamiento es válido cuando no puede darse que con premisas verdaderas la conclusión sea falsa
\end{tcolorbox}

\begin{tcbraster}[raster columns=2, raster equal height=rows, fontupper=\large]
\begin{tcolorbox}[colback=blue!5!blue!20!white,colframe=blue!75!black,title=Razonamientos válidos, subtitle style={boxrule=0.4pt, colback=blue!30!white, colupper=black!75!gray, fontupper=\normalsize}, fontupper=\large, halign=center]
En este contexto, $p$ es la premisa y $q$ la conclusión.
\tcbsubtitle{Modus ponens}
Si $p, \rightarrow q$ \\[8 pt]
$ \dfrac{p}{q} $
\tcbsubtitle{Modus tollens}
Si $p, \rightarrow q$ \\[8 pt]
$ \dfrac{-q}{-p} $
\end{tcolorbox}
\begin{tcolorbox}[colback=blue!5!blue!20!white,colframe=blue!75!black,title=Razonamientos inválidos, subtitle style={boxrule=0.4pt, colback=blue!30!white, colupper=black!75!gray, fontupper=\normalsize}, fontupper=\large, halign=center]
	\tcbsubtitle{Afirmación del consecuente}
	Si $p, \rightarrow q$ \\[8 pt]
	$ \dfrac{q}{p} $
	\tcbsubtitle{Negación del antecedente}
	Si $p, \rightarrow q$ \\[8pt]
	$ \dfrac{-p}{-q} $
\end{tcolorbox}
\end{tcbraster}

\subsubsectionx{4. Términos}

Hay varios tipos de términos a considerar.

\begin{itemize}
	\item Teóricos-observacionales: Los teóricos pertenecen a una jerga específica de una ciencia y dominio, en tanto los observacionales están vinculados a lo que se ve. Ambos pueden variar en el tiempo debido a avances. 
	\item Preteóricos: términos ya aceptados por varias ciencias, pero aún no consolidados en el contexto de una teoría.
	\item Operacionalistas-instrumentalistas: incluye términos clasificatorios, comparativos, métricos. Su significado también puede mutar.
\end{itemize}

\subsubsectionx{5. Enunciados}

\begin{itemize}
	\item Enunciados analíticos: funcionan como axiomas.
	\begin{itemize}
		\item Lógicos.
		\item Matemáticos.
		\item Por sinonimia (semántica).
		\item Definicionales.
	\end{itemize}
	\item Enunciados sintéticos
	\begin{itemize}
		\item Empíricos singulares: hechos o situaciones particulares.
		\item Empíricos generales: ``universales''.
		\item Teóricos: inferencias.
		\item No fácticos: de tipo filosófico. Es mucho más complejo establecer verdad o falsedad.
	\end{itemize}
\end{itemize}

\subsectionx{III. Tipos de conocimiento}

\subsubsectionx{1. Natural}

El conocimiento natural está asociado al lenguaje, al sentido común, el sentido espacial y la socialización. Según F. Bacon, este tipo de conocimiento puede presentar distorsiones:

\begin{itemize}
	\item Tribu: aceptar lo establecido, lo común a todos. Incluye ``wishful thinking''.
	\item Spews: distorsiones a nivel individual.
	\item Fori: distorsiones del lenguaje.
	\item Theatri: dogmas tradicionales.
\end{itemize}

Por lo tanto, este tipo de conocimiento puede acarrear estereotipos y prejuicios. Esto choca con el pensamiento crítico esencial para la ciencia.

\subsubsectionx{2. Conocimiento científico}

Podría argumentarse que una versión inicial de lo que hoy se entiende como conocimiento científico surgió en la antigua Grecia, de la mano de pensadores como Tales de Mileto y Pitágoras. Floreció en el período helenístico, combinando teorías griegas con técnicas egipcias. La ciudad de Alejandría fue el epicentro de estos avances.

Características del conocimiento científico:

\begin{itemize}
	\item Busca unir la teoría y la experimentación.
	\item Es una \textbf{búsqueda constante de la verdad}.
	\item Racional, sostenido por argumentos.
	\item Crítico, como opuesto a lo dogmático. Debe tener autocorrectividad.
	\item Sistematicidad e integración de diversas áreas.
	\item Se busca la precisión en el lenguaje, evitando la ambigüedad.
	\item Métodos propios y especiales.
	\item Se buscan leyes generales sin perder la singularidad.	
\end{itemize}
	
\end{document}
